\section*{第 \ref{ch:g10:exercise-and-energy} 章——运动与人体的能量需求}

\begin{solution}{exo:g10:exercise-and-energy:1}
人在完全静止时为维持生命所消耗的能量——心脏、
大脑、体温及\omterm{def:g10:cells-common-unit:cell}{细胞}维持所需：约 \qty{80}{W}，
每天约 \qty{7000}{kJ}。
\end{solution}

\begin{solution}{exo:g10:exercise-and-energy:2}
每分钟 $2.0 \times 20 = \qty{40}{kJ}$，即 $40000/60 \approx
\qty{670}{W}$。
\end{solution}

\begin{solution}{exo:g10:exercise-and-energy:3}
肌肉的磷酸原储备（约十秒）；乳酸\omterm{def:g10:cell-metabolism:fermentation}{发酵}
（全力运动一至两分钟）；\omterm{prop:g10:cell-metabolism:respiration}{细胞呼吸}（数小时）。
\end{solution}

\begin{solution}{exo:g10:exercise-and-energy:4}
人体能够摄取并利用氧气的最高速率。
$3500/70 = \qty{50}{mL/min}$，即每千克体质量。
\end{solution}

\begin{solution}{exo:g10:exercise-and-energy:5}
肌肉，约 \qty{400}{g}，供储存它的肌肉自身使用；
肝脏，约 \qty{100}{g}，以葡萄糖形式释入血液，供
大脑及其他组织使用。
\end{solution}

\begin{solution}{exo:g10:exercise-and-energy:6}
\qty{2.7}{L/min}：$2.7 \times 20 = \qty{54}{kJ}$ 每分钟，即
\qty{900}{W}。热量：$900 - 200 = \qty{700}{W}$。
\end{solution}

\begin{solution}{exo:g10:exercise-and-energy:7}
功 $mgh = 60 \times 9.8 \times 800 \approx \qty{470}{kJ}$。在 25\%
效率下，肌肉约消耗 \qty{1900}{kJ}。以 \qty{7200}{s} 计，
约比静止时高 \qty{260}{W}。
\end{solution}

\begin{solution}{exo:g10:exercise-and-energy:8}
100 米短跑几乎完全依靠磷酸原储备供能：没有
氧债，乳酸很少。400 米比赛约持续
一分钟，其能量一半来自乳酸\omterm{def:g10:cell-metabolism:fermentation}{发酵}：乳酸
累积（即酸痛感），而需要清除乳酸的呼吸作用
使跑者此后数分钟仍大口喘息。
\end{solution}

\begin{solution}{exo:g10:exercise-and-energy:9}
\omterm{def:g10:chemistry-of-life:families}{糖类}提供的能量 $0.6 \times 3000 = \qty{1800}{kJ}$；\omterm{def:g10:exercise-and-energy:glycogen}{糖原}
质量 $1800/17 \approx \qty{106}{g}$；释放的水约 \qty{320}{g}。
\end{solution}

\begin{solution}{exo:g10:exercise-and-energy:10}
\omterm{def:g10:cell-metabolism:metabolism}{代谢}功率约 \qty{600}{W}：$\num{1.1e6}/600 \approx
\qty{1830}{s}$，约 30 分钟。
\end{solution}

\begin{solution}{exo:g10:exercise-and-energy:11}
肌肉释放的能量有四分之三转化为热量，剧烈运动时可达
数百瓦。出汗可排出热量：蒸发
每克水约需 \qty{2.4}{kJ}，同时血液流向皮肤，
将热量带到那里。若不出汗，体温将
每隔几分钟升高一度。
\end{solution}

\begin{solution}{exo:g10:exercise-and-energy:12}
上坡时，所需功率与体质量成正比：
\qty{60}{kg} 的受试者每千克体质量可达 \qty{60}{mL/min}，
相比 \num{40}，攀爬更快。在实验室自行车上，功率
与质量无关，两人均为 \qty{3.6}{L/min}，
两人都能维持相同的瓦数。
\end{solution}

\begin{solution}{exo:g10:exercise-and-energy:13}
缺少的 \qty{0.4}{L/min} 氧当量由
肌肉中的乳酸\omterm{def:g10:cell-metabolism:fermentation}{发酵}提供。乳酸累积，几分钟内
肌肉便会灼痛，跑者被迫减速至
仅靠呼吸作用即可维持的速度。
\end{solution}

\begin{solution}{exo:g10:exercise-and-energy:14}
短时间最大强度的爆发依靠磷酸原储备和\omterm{def:g10:cell-metabolism:fermentation}{发酵}供能，
两者都完全不消耗脂肪；脂肪只由呼吸作用消耗，而呼吸作用
需要数分钟才能达到全速，只在长时间的
中等强度运动中占主导。几次爆发的总能量也很少，
与持续一小时的稳定运动相比。
\end{solution}

\begin{solution}{exo:g10:exercise-and-energy:15}
能量缺口 $25000 - 8500 - 6000 = \qty{10500}{kJ}$，即约
\qty{280}{g} 脂肪。赛程中进食，可维持血糖供大脑使用，
并节省\omterm{def:g10:exercise-and-energy:glycogen}{糖原}用于爬坡，爬坡时
所需功率远超脂肪单独所能供应的量；脂肪则负担
全天较平稳的部分。
\end{solution}

\begin{solution}{pb:g10:exercise-and-energy:1}
\textbf{1.} $4.0 \times 60 \times 42.2 \approx \qty{10100}{kJ}$。

\textbf{2.} $\qty{3}{h}\,\qty{30}{min} = \qty{12600}{s}$：
$\num{1.01e7}/12600 \approx \qty{800}{W}$。

\textbf{3.} $10100/20 \approx \qty{506}{L}$；在 \qty{210}{min} 内，约
\qty{2.4}{L/min}。

\textbf{4.} $2400/60 = \qty{40}{mL/min}$，即每千克体质量 $40/52
\approx 77\%$，约占她最大值的比例。

\textbf{5.} $(2.4 - 0.2)/2.4 \approx 92\%$。

\textbf{6.} $380 \times 17 = \qty{6460}{kJ}$。

\textbf{7.} 比赛每千米耗能 \qty{240}{kJ}：$6460/240
\approx \qty{27}{km}$。

\textbf{8.} \omterm{def:g10:chemistry-of-life:families}{糖类}消耗 $0.8 \times 240 = \qty{192}{kJ/km}$：
$6460/192 \approx \qty{34}{km}$——撞墙点。

\textbf{9.} 脂肪供能 $0.2 \times 10100 \approx \qty{2020}{kJ}$，即
$2020/38 \approx \qty{53}{g}$：约占她 \qty{9}{kg} 的 0.6\%。

\textbf{10.} 脂肪只能通过呼吸作用消耗，且速率有限：
它只能提供配速所需功率的约一半。若无
\omterm{def:g10:exercise-and-energy:glycogen}{糖原}，肌肉便无法足够快地生成 ATP，配速随之下降。

\textbf{11.} $60 \times 3.5 = \qty{210}{g}$，即 $210 \times 17 \approx
\qty{3570}{kJ}$，可覆盖 $3570/192 \approx \qty{19}{km}$ 的
\omterm{def:g10:chemistry-of-life:families}{糖类}消耗。

\textbf{12.} 可用\omterm{def:g10:chemistry-of-life:families}{糖类} $6460 + 3570 = \qty{10030}{kJ}$：
$10030/192 \approx \qty{52}{km}$，超过终点。足够。

\textbf{13.} 增加\omterm{def:g10:exercise-and-energy:glycogen}{糖原} $+\qty{200}{g}$，水 $+ \qty{600}{g}$：
\qty{0.8}{kg}，约增加 $0.8 \times 4.0 \approx \qty{3}{kJ}$ 每千米的
额外耗能——相对 240 可忽略不计。

\textbf{14.} \omterm{def:g10:exercise-and-energy:glycogen}{糖原} $580 \times 17 = \qty{9860}{kJ}$，加上饮料：
共约 \qty{13400}{kJ}；\omterm{def:g10:chemistry-of-life:families}{糖类}消耗 $0.55 \times 240 =
\qty{132}{kJ/km}$，可维持超过 \qty{100}{km}。撞墙点已远远后移，
超过终点，余量充足。

\textbf{15.} 脑只消耗葡萄糖，这些葡萄糖来自
肝脏\omterm{def:g10:exercise-and-energy:glycogen}{糖原}。储备耗尽时血糖下降，脑
以头晕和极度疲劳发出紧急信号，
无论肌肉中还剩多少\omterm{def:g10:exercise-and-energy:glycogen}{糖原}。

\textbf{16.} $3.6 \times 55 \times 42.2 \approx \qty{8360}{kJ}$；
$8360/20 = \qty{418}{L}$，在 \qty{125}{min} 内，约 \qty{3.3}{L/min}。

\textbf{17.} $3340/55 \approx \qty{61}{mL/min}$，即每千克体质量
$61/80 \approx 76\%$——与娜迪亚相同。上限
不同，而所用比例相同。

\textbf{18.} 跑步经济性指每千克体质量、每千米的耗能。对
\qty{55}{kg} 的跑者，$(4.0 - 3.6) \times 55 \times 42.2 \approx
\qty{930}{kJ}$，全程节省约 9\%。

\textbf{19.} \omterm{def:g10:exercise-and-energy:glycogen}{糖原} \qty{8500}{kJ}；所需\omterm{def:g10:chemistry-of-life:families}{糖类} $0.7 \times
8360 \approx \qty{5850}{kJ}$：储备足够，
无需饮用饮料。

\textbf{20.} 若不做准备，娜迪亚的\omterm{def:g10:exercise-and-energy:glycogen}{糖原}将在
第 34 千米处耗尽；比赛前的\omterm{def:g10:exercise-and-energy:glycogen}{糖原}负荷、比赛中饮用
\omterm{def:g10:chemistry-of-life:families}{糖类}饮料，以及提高脂肪供能占比的训练，
均可使撞墙点后移至终点线之外。
\end{solution}
